\SourceLeaf{75}
\Needspace{10\baselineskip}\begin{center}
\textbf{REMARKS ON THE TWO PRECEDING FRAGMENTS ON A THEOREM OF EULER.}
\end{center}

\noindent\hspace*{\parindent}\textsc{Gauss} treats this theorem of \textsc{Euler}: “The function
$V(x, y, p, q, r, s, .\,.\,.)$ of $x, y, p=\dfrac{dy}{dx}$, $q=\dfrac{dp}{dx}$, $r=\dfrac{dq}{dx}$, $s=\dfrac{dr}{dx}$, $.\,.\,.$
has the necessary and sufficient condition
\[
N-\dfrac{dP}{dx}+\dfrac{d^2Q}{dx^2}-\dfrac{d^3R}{dx^3}+\dfrac{d^4S}{dx^4}-.\,.\,.=0
\]
for yielding an exact differential $Vdx$, where $N, P, Q, R, S, .\,.\,.$ denote the partial derivatives
\[
\begin{aligned}N=\dfrac{\partial V}{\partial y},\quad P=\dfrac{\partial V}{\partial p},\quad Q=\dfrac{\partial V}{\partial q},\\ R=\dfrac{\partial V}{\partial r},\quad S=\dfrac{\partial V}{\partial s},\ .\,.\,.
\end{aligned}
\]
of $V$.” \textsc{Euler} stated this theorem and its converse in the third part of his
‘Institutiones calculi integralis’, in \S\,92 of the appendix on the calculus of variations, and proved
them from the preceding developments of the calculus of variations.
\textsc{Gauss} recorded his direct proof of the theorem, reproduced above, at the relevant place
in his personal copy of \textsc{Euler}’s Institutiones calculi integralis. \textsc{Gauss} wrote
the proof of the converse developed in the second fragment in his copy of
\textsc{Klügel}’s Lexicon, Vol.\,I.

\begin{flushright}
\textsc{Fricke}.
\end{flushright}

